\section{Conclusi\'on}
\label{sec:conclusion}

Esta tesis ha presentado MalWhere, un pipeline que combina
an\'alisis est\'atico y din\'amico de malware para producir
inteligencia de amenazas con confianza cuantificada y conforme a
est\'andares: mapeos de t\'ecnicas MITRE ATT\&CK exportados como
paquetes STIX~2.1 y eventos MISP en vivo. Dos mecanismos lo extienden
m\'as all\'a de un dise\~no est\'atico-m\'as-din\'amico convencional: un
modelo de reconciliaci\'on de confianza que reconoce concordancia tanto
entre fuentes como entre niveles de especificidad de t\'ecnica
(\S\ref{sec:reconciliation}), y un bucle de reenv\'io que atribuye
correctamente la capacidad distribuida de una muestra multietapa al
componente que realmente la presenta, en lugar de colapsarla en la
puntuaci\'on propia del binario padre (\S\ref{sec:resubmission}).
Frente a las herramientas revisadas en \S\ref{sec:related-work}, el
pipeline no pretende sustituir a un agregador de reputaci\'on amplio
como VirusTotal, al que ning\'un pipeline de investigaci\'on de un solo
equipo puede igualar en escala bruta de corpus de env\'ios; es el paso
que viene despu\'es, convirtiendo esa primera se\~nal en un mapeo
reconciliado, por t\'ecnica y ligado a evidencia para la muestra
concreta que realmente lo necesita.

Validado frente a tres familias de malware con arquitecturas
sustancialmente distintas (un binario de ransomware de una sola etapa,
un infostealer .NET rico en funcionalidades, y una cadena de cuatro
componentes cargador/rootkit/RAT) contra verdad de referencia
construida a partir de ingenier\'ia inversa manual a nivel de
funci\'on, el pipeline alcanz\'o puntuaciones F1 a nivel de familia de
0.96, 0.89 y 0.84 respectivamente, con precisi\'on alta
en todos los casos (\S\ref{sec:evaluation}). Esas cifras solo son tan
fiables como el proceso que las produjo, raz\'on por la cual esta tesis
otorga igual peso al rastro de auditor\'ia que las respalda: 15
discrepancias entre el resultado del pipeline y la verdad de
referencia se rastrearon individualmente hasta una causa concreta y
verificable, no se aceptaron ni descartaron solo por la cifra agregada
(\S\ref{sec:fp-audit}); se encontr\'o y corrigi\'o un error real de
agregaci\'on que hab\'ia sustituido silenciosamente el binario analizado
incorrecto (\S\ref{sec:case-installer}); y la cobertura de detecci\'on
est\'atica se ampli\'o de unas 30 a aproximadamente 85 identificadores
de t\'ecnica MITRE, cada adici\'on obtenida de las propias
descripciones de t\'ecnica de MITRE y confirmada para no introducir
nuevos hallazgos en el conjunto validado antes de confiar en ella
(\S\ref{sec:coverage}).

El argumento central de este trabajo es que este rastro de
auditor\'ia, cada regla ligada a evidencia verificable, cada
discrepancia diagnosticada a una causa concreta, cada extensi\'on
comprobada frente a regresiones antes de conservarse, es una
contribuci\'on m\'as defendible que las cifras F1 destacadas por s\'i
solas, y es la propiedad con mayor probabilidad de transferirse cuando
el pipeline se encuentre con una familia de malware para la que no ha
sido afinada. Precisamente porque esa transferencia todav\'ia no se ha
medido, ampliar el conjunto de familias validadas
(\S\ref{sec:future-work}) se identifica como el siguiente paso de
mayor valor, junto con la extensi\'on concretamente acotada al
an\'alisis de botnets IoT basadas en Linux/ELF una vez que sus dos
prerrequisitos faltantes, reglas de detecci\'on realmente conscientes
de ELF y una muestra Linux con verdad de referencia, est\'en
disponibles.
